\documentclass[10pt,a4paper]{amsart}
\usepackage{amsmath,amssymb,amsthm,graphicx}
\usepackage{verbatim}
\newenvironment{reference}{\comment}{\endcomment}
%\newenvironment{reference}{}{}
\newenvironment{slogan}{\comment}{\endcomment}
\newenvironment{history}{\comment}{\endcomment}
% For commutative diagrams we use Xy-pic
\usepackage[all]{xy}

% We use 2cell for 2-commutative diagrams.
\xyoption{2cell}
% Theorem environments.
%
\theoremstyle{plain}
\newtheorem{theorem}[subsection]{Theorem}
\newtheorem{proposition}[subsection]{Proposition}
\newtheorem{lemma}[subsection]{Lemma}

\theoremstyle{definition}
\newtheorem{definition}[subsection]{Definition}
\newtheorem{example}[subsection]{Example}
\newtheorem{exercise}[subsection]{Exercise}
\newtheorem{situation}[subsection]{Situation}

\theoremstyle{remark}
\newtheorem{remark}[subsection]{Remark}
\newtheorem{remarks}[subsection]{Remarks}

\numberwithin{equation}{subsection}

% Macros
%
\def\lim{\mathop{\mathrm{lim}}\nolimits}
\def\colim{\mathop{\mathrm{colim}}\nolimits}
\def\Spec{\mathop{\mathrm{Spec}}}
\def\Hom{\mathop{\mathrm{Hom}}\nolimits}
\def\Ext{\mathop{\mathrm{Ext}}\nolimits}
\def\SheafHom{\mathop{\mathcal{H}\!\mathit{om}}\nolimits}
\def\SheafExt{\mathop{\mathcal{E}\!\mathit{xt}}\nolimits}
\def\Sch{\mathit{Sch}}
\def\Mor{\mathop{\mathrm{Mor}}\nolimits}
\def\Ob{\mathop{\mathrm{Ob}}\nolimits}
\def\Sh{\mathop{\mathit{Sh}}\nolimits}
\def\NL{\mathop{N\!L}\nolimits}
\def\CH{\mathop{\mathrm{CH}}\nolimits}
\def\proetale{{pro\text{-}\acute{e}tale}}
\def\etale{{\acute{e}tale}}
\def\QCoh{\mathit{QCoh}}
\def\Ker{\mathop{\mathrm{Ker}}}
\def\Im{\mathop{\mathrm{Im}}}
\def\Coker{\mathop{\mathrm{Coker}}}
\def\Coim{\mathop{\mathrm{Coim}}}

% Boxtimes
%
\DeclareMathSymbol{\boxtimes}{\mathbin}{AMSa}{"02}

%
% Macros for moduli stacks/spaces
%
\def\QCohstack{\mathcal{QC}\!\mathit{oh}}
\def\Cohstack{\mathcal{C}\!\mathit{oh}}
\def\Spacesstack{\mathcal{S}\!\mathit{paces}}
\def\Quotfunctor{\mathrm{Quot}}
\def\Hilbfunctor{\mathrm{Hilb}}
\def\Curvesstack{\mathcal{C}\!\mathit{urves}}
\def\Polarizedstack{\mathcal{P}\!\mathit{olarized}}
\def\Complexesstack{\mathcal{C}\!\mathit{omplexes}}
% \Pic is the operator that assigns to X its picard group, usage \Pic(X)
% \Picardstack_{X/B} denotes the Picard stack of X over B
% \Picardfunctor_{X/B} denotes the Picard functor of X over B
\def\Pic{\mathop{\mathrm{Pic}}\nolimits}
\def\Picardstack{\mathcal{P}\!\mathit{ic}}
\def\Picardfunctor{\mathrm{Pic}}
\def\Deformationcategory{\mathcal{D}\!\mathit{ef}}
\title{Being separated, quasi-separated, quasi-DM or DM is inherited by the fibre products of test spaces and by the Isom sheaves}
\author{The Stacks Project Authors}
\date{stacks-project:stacks-morphisms.tex:local-snapshot-20261012; GNU FDL 1.2 or later}
\begin{document}
\maketitle
\noindent\textit{Editorial note.} All delivered mathematical text is the author's own, copied from the pinned Stacks Project chapter source sources/stacks-algebraic-stacks-morphisms/pinned-original-source.tex; no statement or proof step was rewritten, reordered or removed.\par\medskip
\noindent\textit{Editorial note.} Preamble provenance: the Stacks Project keeps one single shared preamble for every chapter; this card takes it directly from the staged sources/stacks-algebraic-stacks-morphisms/preamble.tex as three verbatim spans (lines 8-12, 14-18 and 154-217), so no derived file was written and nothing was added or redefined. The wrapper class and the base packages are supplied by the assembler.\par\medskip
\noindent\textit{Editorial note.} Bibliography provenance: neither the counted statement nor its proof contains a citation (this lane rescans the whole span before writing the card), so the card carries no bibliography and no entry was reconstructed.\par\medskip
\noindent\textit{Editorial note.} Selection provenance: this card was opened under the file-based control contract (CONTROL.md), with the mechanical pre-card gate run on its statement and proof interval before the card was written.\par\medskip
\noindent\textit{Allowed standard prerequisites (editorial).}\begin{itemize}
\item The characterisation of separated, quasi-separated, quasi-DM and DM morphisms of algebraic stacks by the corresponding property of the diagonal. Invoked inside the counted proof (source lines 1885-1889). Converts the property of the morphism into a property of the diagonal, which is what the displayed cartesian square tests.
\item The description of the Isom sheaf of two objects over a common test space as a fibre product of the double fibre product, and stability of the properties under base change. Used inside the counted proof (source lines 1892-1897). Exhibits the morphism of the second assertion as a base change of the morphism of the first.
\end{itemize}\medskip
\begin{lemma}
\label{lemma-separated-implies-isom}
Let $f : \mathcal{X} \to \mathcal{Y}$ be a separated
(resp.\ quasi-separated, resp.\ quasi-DM, resp.\ DM)
morphism of algebraic stacks. Then
\begin{enumerate}
\item given algebraic spaces $T_i$, $i = 1, 2$ and morphisms
$x_i : T_i \to \mathcal{X}$, with $y_i = f \circ x_i$ the morphism
$$
T_1 \times_{x_1, \mathcal{X}, x_2} T_2 \longrightarrow
T_1 \times_{y_1, \mathcal{Y}, y_2} T_2
$$
is proper (resp.\ quasi-compact and quasi-separated,
resp.\ locally quasi-finite, resp.\ unramified),
\item given an algebraic space $T$ and morphisms
$x_i : T \to \mathcal{X}$, $i = 1, 2$, with $y_i = f \circ x_i$ the morphism
$$
\mathit{Isom}_\mathcal{X}(x_1, x_2) \longrightarrow
\mathit{Isom}_\mathcal{Y}(y_1, y_2)
$$
is proper (resp.\ quasi-compact and quasi-separated,
resp.\ locally quasi-finite, resp.\ unramified).
\end{enumerate}
\end{lemma}

\begin{proof}
Proof of (1). Observe that the diagram
$$
\xymatrix{
T_1 \times_{x_1, \mathcal{X}, x_2} T_2 \ar[d] \ar[r] &
T_1 \times_{y_1, \mathcal{Y}, y_2} T_2 \ar[d] \\
\mathcal{X} \ar[r] & \mathcal{X} \times_\mathcal{Y} \mathcal{X}
}
$$
is cartesian. Hence this follows from the fact that $f$ is separated
(resp.\ quasi-separated, resp.\ quasi-DM, resp.\ DM)
if and only if the diagonal is proper
(resp.\ quasi-compact and quasi-separated,
resp.\ locally quasi-finite, resp.\ unramified).

\medskip\noindent
Proof of (2). This is true because
$$
\mathit{Isom}_\mathcal{X}(x_1, x_2) =
(T \times_{x_1, \mathcal{X}, x_2} T) \times_{T \times T, \Delta_T} T
$$
hence the morphism in (2) is a base change of the morphism in (1).
\end{proof}

\end{document}
